\clearpage
\SourceLeaf{381}%
\begin{center}\large\bfseries STAND MEINER UNTERSUCHUNG\\ ÜBER DIE UMFORMUNG DER FLÄCHEN.\end{center}
\begin{center}Dec. 13. 1822.\end{center}
\begin{center}[1.]\end{center}

[Es seien
\begin{description}
\item[$x,y,z$] die rechtwinkligen Coordinaten eines Punktes der Fläche,
\item[$A,B,C$] die Cosinus der Winkel, welche die Normale der Fläche in diesem Punkte mit den Coordinatenaxen bildet. Dann ist:]
\end{description}
\begin{equation}\tag*{[1]} A\,dx+B\,dy+C\,dz=0\end{equation}
\begin{equation}\tag*{[2]} AA+BB+CC=1.\end{equation}

[Die Coordinaten $x,y,z$ mögen durch zwei unabhängige Veränderliche $t,u$ ausgedrückt werden. Man setze:]
\begin{equation}\tag*{[3]}
\left\{\begin{array}{lll}
A'=\dfrac{\partial A}{\partial t},&B'=\dfrac{\partial B}{\partial t},&C'=\dfrac{\partial C}{\partial t},\\[6pt]
A''=\dfrac{\partial A}{\partial u},&B''=\dfrac{\partial B}{\partial u},&C''=\dfrac{\partial C}{\partial u};
\end{array}\right.
\end{equation}
\begin{equation}\tag*{[4]}
\left\{\begin{array}{lll}
a=\dfrac{\partial x}{\partial t},&b=\dfrac{\partial y}{\partial t},&c=\dfrac{\partial z}{\partial t},\\[6pt]
\alpha=\dfrac{\partial x}{\partial u},&\beta=\dfrac{\partial y}{\partial u},&\gamma=\dfrac{\partial z}{\partial u};
\end{array}\right.
\end{equation}
